\documentclass[tikz,border=10pt]{standalone}
\usepackage[no-math]{fontspec}
\usepackage[brazil]{babel}
\usepackage{amsmath,amssymb}
\IfFileExists{C:/Windows/Fonts/Lato_Regular.ttf}{%
  \setmainfont{Lato}[
    Path = C:/Windows/Fonts/,
    UprightFont = Lato_Regular.ttf,
    BoldFont = Lato_Bold.ttf,
    ItalicFont = Lato_Italic.ttf,
    BoldItalicFont = Lato_BoldItalic.ttf
  ]
  \setsansfont{Lato}[
    Path = C:/Windows/Fonts/,
    UprightFont = Lato_Regular.ttf,
    BoldFont = Lato_Bold.ttf,
    ItalicFont = Lato_Italic.ttf,
    BoldItalicFont = Lato_BoldItalic.ttf
  ]
}{%
  \RequirePackage[default]{lato}
}
\renewcommand{\familydefault}{\sfdefault}

\usepackage{xcolor}
\usepackage{tikz}
\usepackage{pgfplots}
\pgfplotsset{compat=1.18}
\usetikzlibrary{arrows.meta,positioning,shapes.geometric,fit,backgrounds,calc,decorations.pathreplacing,patterns}

% Paleta institucional UEFS / Bioengenharia de Solos
\definecolor{bsNavy}{HTML}{1F4268}
\definecolor{bsNavyDark}{HTML}{152E4A}
\definecolor{bsGold}{HTML}{B38A3D}
\definecolor{bsGoldDark}{HTML}{8C6927}
\definecolor{bsTeal}{HTML}{0E7C7B}
\definecolor{bsForest}{HTML}{2D6A4F}
\definecolor{bsAlert}{HTML}{9E2A2B}
\definecolor{bsEarth}{HTML}{8C5A3C}
\definecolor{bsWater}{HTML}{2E86AB}
\definecolor{bsInk}{HTML}{26313D}
\definecolor{bsMuted}{HTML}{5B6573}

\definecolor{bgNavy}{HTML}{EDF2F7}
\definecolor{bgGold}{HTML}{F8F3E9}
\definecolor{bgTeal}{HTML}{EBF6F5}
\definecolor{bgForest}{HTML}{EEF5F0}
\definecolor{bgAlert}{HTML}{FDF2F2}

\tikzset{
  >=Stealth,
  bsBox/.style={
    draw=#1, fill=white, rounded corners=5pt, line width=1.15pt,
    inner sep=7pt, align=center, font=\footnotesize\sffamily, text=bsInk
  },
  bsTint/.style={
    draw=#1, fill=#1!8!white, rounded corners=5pt, line width=1.15pt,
    inner sep=7pt, align=center, font=\footnotesize\sffamily, text=bsInk
  },
  bsBadge/.style={
    fill=#1, text=white, rounded corners=9pt, inner xsep=8pt, inner ysep=3pt,
    font=\scriptsize\bfseries\sffamily
  },
  bsTitle/.style={font=\normalsize\bfseries\sffamily, text=bsNavy, align=center},
  bsSub/.style={font=\footnotesize\sffamily, text=bsMuted, align=center},
  bsArrow/.style={->, >=Stealth, line width=1.35pt},
  bsDArrow/.style={<->, >=Stealth, line width=1.35pt}
}


\begin{document}
\begin{tikzpicture}[font=\sffamily]

% ============================================================
% TÍTULO GERAL
% ============================================================
\node[bsTitle, anchor=west] at (0, 8.65)
  {Parede Krainer Vegetada: Geometria da Grelha e Verificação como Muro de Gravidade};
\node[bsSub, anchor=west] at (0, 8.20)
  {Estrutura em fogueira de toras com solo compactado, ramos vivos alternados e critérios de estabilidade};

% ============================================================
% PAINEL ESQUERDO: SEÇÃO TRANSVERSAL DA PAREDE KRAINER
% ============================================================
\begin{scope}[shift={(0,0)}]
  \draw[bsNavy!30, rounded corners=6pt, line width=1.0pt, fill=white]
    (0,-0.75) rectangle (9.65, 7.70);
  \node[bsBadge=bsEarth, rounded corners=5pt, anchor=north west] at (0.30, 7.45) {GEOMETRIA DA GRELHA};

  % Solo natural atrás e base escavada abaixo do solapamento
  \fill[bsEarth!18]
    (0.45, 0.75) -- (0.45, 5.95) -- (2.95, 5.95) -- (3.55, 1.35) -- (8.85, 1.35) -- (8.85, 0.75) -- cycle;

  % Linha de referência vertical para mostrar inclinação de 10-15% para dentro
  \draw[bsMuted, dotted, line width=0.9pt] (5.95, 1.35) -- (5.95, 6.20);

  % Solo compactado interno da grelha (inclinado para dentro ~12%)
  \fill[bsGold!28]
    (3.55, 1.35) -- (5.95, 1.35) -- (5.35, 5.45) -- (2.95, 5.45) -- cycle;

  % Toras transversais e longitudinais em 5 níveis
  \foreach \i in {0,1,2,3,4}{
    \pgfmathsetmacro{\y}{1.45 + \i*0.92}
    \pgfmathsetmacro{\xL}{3.55 - \i*0.13}
    \pgfmathsetmacro{\xR}{5.95 - \i*0.13}

    \filldraw[fill=bsEarth!55, draw=bsEarth!90!black, line width=0.8pt]
      (\xL-0.06, \y+0.12) -- (\xR+0.06, \y+0.12) -- (\xR+0.04, \y+0.32) -- (\xL-0.08, \y+0.32) -- cycle;
    \filldraw[fill=bsGoldDark!75, draw=bsEarth!90!black, line width=0.85pt]
      (\xL+0.10, \y) circle (0.14);
    \filldraw[fill=bsGoldDark!75, draw=bsEarth!90!black, line width=0.85pt]
      (\xR-0.10, \y) circle (0.14);
  }

  % Camadas alternadas de ramos vivos atravessando o muro e brotando na face
  \foreach \y/\xF in {2.02/5.88, 2.94/5.75, 3.86/5.62, 4.78/5.49}{
    \draw[bsEarth!85!black, line width=2.1pt, line cap=round] (\xF-3.10, \y+0.25) -- (\xF+0.32, \y-0.05);
    \draw[bsForest, line width=1.5pt, line cap=round] (\xF+0.15, \y) -- (\xF+0.78, \y+0.36);
    \draw[bsForest, line width=1.35pt, line cap=round] (\xF+0.22, \y-0.03) -- (\xF+0.86, \y+0.10);
    \fill[bsForest!85!white] (\xF+0.78, \y+0.36) ellipse (0.18 and 0.09);
    \fill[bsForest!75!white] (\xF+0.86, \y+0.10) ellipse (0.17 and 0.08);
  }

  % Cota de solapamento na base
  \draw[bsAlert, dashed, line width=0.95pt] (3.30, 1.82) -- (7.65, 1.82);
  \node[font=\tiny\bfseries\sffamily, text=bsAlert, anchor=south east] at (8.80, 1.86) {Cota de solapamento};
  \draw[bsEarth, line width=1.1pt] (5.95, 1.82) -- (8.85, 1.82);

  % Indicação de inclinação 10-15% (acima do topo, sem cruzamentos)
  \draw[bsDArrow, bsNavy, line width=0.95pt] (5.42, 5.78) -- (5.95, 5.78);
  \node[font=\scriptsize\bfseries\sffamily, text=bsNavy, anchor=south] at (5.68, 5.84) {10–15\%};

  % Rótulos e setas (totalmente desimpedidos)
  \node[font=\scriptsize\bfseries\sffamily, text=bsEarth!90!black, align=left, anchor=west] at (0.55, 6.32)
    {Solo natural\\e enraizamento};

  \node[font=\scriptsize\bfseries\sffamily, text=bsGoldDark, align=left, anchor=west] (lToras) at (6.50, 6.25)
    {Toras de madeira\\tratada / nativa};
  \draw[bsArrow, bsGoldDark, line width=0.95pt] (lToras.south west) -- (5.58, 5.20);

  \node[font=\scriptsize\bfseries\sffamily, text=bsForest, align=left, anchor=west] (lCam) at (6.55, 4.15)
    {Camadas alternadas\\de solo e ramos vivos};
  \draw[bsArrow, bsForest, line width=0.95pt] (lCam.west) -- (6.30, 4.02);

  \node[font=\scriptsize\bfseries\sffamily, text=bsNavyDark, anchor=north] at (4.75, 0.56)
    {Base engastada (muro de gravidade)};
\end{scope}

% ============================================================
% PAINEL DIREITO: VERIFICAÇÕES GEOTÉCNICAS CLÁSSICAS
% ============================================================
\begin{scope}[shift={(10.05,0)}]
  \draw[bsAlert, rounded corners=6pt, line width=1.2pt, fill=bgAlert]
    (0,-0.75) rectangle (5.85, 7.70);
  \node[bsBadge=bsAlert, rounded corners=5pt, anchor=north west] at (0.30, 7.45) {ESTABILIDADE GLOBAL};

  \node[font=\small\bfseries\sffamily, text=bsNavyDark, anchor=north west, text width=5.15cm]
    at (0.30, 6.86) {Verificações Geotécnicas\\Clássicas};

  \node[bsBox=bsNavy, anchor=north west, text width=4.95cm, align=left] (v1) at (0.30, 5.85)
    {\textbf{\color{bsNavy}1. Tombamento}\\
     $FS_{\text{tombamento}} \ge 1{,}5$ em relação ao pé frontal da grelha.};

  \node[bsBox=bsNavy, anchor=north west, text width=4.95cm, align=left] (v2) at ($(v1.south west)+(0,-0.22)$)
    {\textbf{\color{bsNavy}2. Deslizamento na Base}\\
     $FS_{\text{deslizamento}} \ge 1{,}5$ sob ação do empuxo ativo do maciço.};

  \node[bsBox=bsAlert, anchor=north west, text width=4.95cm, align=left] (v3) at ($(v2.south west)+(0,-0.22)$)
    {\textbf{\color{bsAlert}3. Tensão Admissível na Base}\\
     Verificar capacidade de suporte da fundação: $\sigma_{\text{máx}} \le \sigma_{\text{adm}}$.};

  \node[bsBox=bsForest, anchor=north west, text width=4.95cm, align=left] (v4) at ($(v3.south west)+(0,-0.22)$)
    {\textbf{\color{bsForest}4. Drenagem Interna}\\
     Face vazada alivia poropressões e evita acúmulo hidrostático.};
\end{scope}

\end{tikzpicture}
\end{document}
